\documentclass{article}
\usepackage{graphicx} %LaTeX package to import graphics
\usepackage{amsmath}% For the equation* environment
\graphicspath{{../images/}} %configuring the graphicx package
\title{My first LaTeX document}
\author{Hubert Farnsworth\thanks{Funded by the Overleaf team.}}
\date{August 2022}
\begin{document}
\maketitle
\begin{abstract}
  This is a simple paragraph at the beginning of the 
  document. A brief introduction about the main subject.
\end{abstract}

We have now added a title, author and date to our first \LaTeX{} document!

\section{Font Style}

Some of the \textbf{greatest}
discoveries in \underline{science} 
were made by \textbf{\textit{accident}}.

Some of the greatest \emph{discoveries} in science 
were made by accident.

\textit{Some of the greatest \emph{discoveries} 
in science were made by accident.}

\textbf{Some of the greatest \emph{discoveries} 
in science were made by accident.}

\section{Image and Figure}

The universe is immense and it seems to be homogeneous, 
on a large scale, everywhere we look.

% The \includegraphcs command is 
% provided (implemented) by the 
% graphicx package
\includegraphics[width=0.25\textwidth]{alipay}  
 
There's a picture of a galaxy above.

\begin{figure}[h]
  \centering
  \includegraphics[width=0.25\textwidth]{wepay}
  \caption{Wepay}
  \label{fig:wepay}
\end{figure}

As you can see in figure \ref{fig:wepay}, the function grows near the origin. This example is on page \pageref{fig:wepay}.

\section{List}

\begin{itemize}
  \item The individual entries are indicated with a black dot, a so-called bullet.
  \item The text in the entries may be of any length.
\end{itemize}

\begin{enumerate}
  \item This is the first entry in our list.
  \item The list numbers increase with each entry we add.
\end{enumerate}


\section{Math mode}

In physics, the mass-energy equivalence is stated 
by the equation $E=mc^2$, discovered in 1905 by Albert Einstein.


\begin{math}
  E=mc^2
\end{math} is typeset in a paragraph using inline math mode---as is $E=mc^2$, and so too is \(E=mc^2\).


The mass-energy equivalence is described by the famous equation
\[ E=mc^2 \] discovered in 1905 by Albert Einstein. 

In natural units ($c = 1$), the formula expresses the identity
\begin{equation}
E=m
\end{equation}


Subscripts in math mode are written as $a_b$ and superscripts are written as $a^b$. These can be combined and nested to write expressions such as

\[ T^{i_1 i_2 \dots i_p}_{j_1 j_2 \dots j_q} = T(x^{i_1},\dots,x^{i_p},e_{j_1},\dots,e_{j_q}) \]
 
We write integrals using $\int$ and fractions using $\frac{a}{b}$. Limits are placed on integrals using superscripts and subscripts:

\[ \int_0^1 \frac{dx}{e^x} =  \frac{e-1}{e} \]

Lower case Greek letters are written as $\omega$ $\delta$ etc. while upper case Greek letters are written as $\Omega$ $\Delta$.

Mathematical operators are prefixed with a backslash as $\sin(\beta)$, $\cos(\alpha)$, $\log(x)$ etc.


\section{First example}

The well-known Pythagorean theorem \(x^2 + y^2 = z^2\) was proved to be invalid for other exponents, meaning the next equation has no integer solutions for \(n>2\):

\[ x^n + y^n = z^n \]

\section{Second example}

This is a simple math expression \(\sqrt{x^2+1}\) inside text. 
And this is also the same: 
\begin{math}
\sqrt{x^2+1}
\end{math}
but by using another command.

This is a simple math expression without numbering
\[\sqrt{x^2+1}\] 
separated from text.

This is also the same:
\begin{displaymath}
\sqrt{x^2+1}
\end{displaymath}

\ldots and this:
\begin{equation*}
\sqrt{x^2+1}
\end{equation*}


\section{New Line}

After our abstract we can begin the first paragraph, then press ``enter'' twice to start the second one.

This line will start a second paragraph.

I will start the third paragraph and then add \\ a manual line break which causes this text to start on a new line but remains part of the same paragraph. Alternatively, I can use the \verb|\newline|\newline command to start a new line, which is also part of the same paragraph.


\section{Table}

\begin{center}
\begin{tabular}{c c c}
  cell1 & cell2 & cell3 \\ 
  cell4 & cell5 & cell6 \\  
  cell7 & cell8 & cell9    
\end{tabular}
\end{center}


\begin{center}
\begin{tabular}{|c|c|c|} 
  \hline
  cell1 & cell2 & cell3 \\ 
  cell4 & cell5 & cell6 \\ 
  cell7 & cell8 & cell9 \\ 
  \hline
\end{tabular}
\end{center}


\begin{center}
\begin{tabular}{||c c c c||} 
  \hline
  Col1 & Col2 & Col2 & Col3 \\ [0.5ex] 
  \hline\hline
  1 & 6 & 87837 & 787 \\ 
  \hline
  2 & 7 & 78 & 5415 \\
  \hline
  3 & 545 & 778 & 7507 \\
  \hline
  4 & 545 & 18744 & 7560 \\
  \hline
  5 & 88 & 788 & 6344 \\ [1ex] 
  \hline
\end{tabular}
\end{center}


Table \ref{table:data} shows how to add a table caption and reference a table.
\begin{table}[h!]
\centering
\begin{tabular}{||c c c c||} 
  \hline
  Col1 & Col2 & Col2 & Col3 \\ [0.5ex] 
  \hline\hline
  1 & 6 & 87837 & 787 \\ 
  2 & 7 & 78 & 5415 \\
  3 & 545 & 778 & 7507 \\
  4 & 545 & 18744 & 7560 \\
  5 & 88 & 788 & 6344 \\ [1ex] 
  \hline
\end{tabular}
\caption{Table to test captions and labels.}
\label{table:data}
\end{table}


\section{Comments}

% This line here is a comment. It will not be typeset in the document.
% https://www.overleaf.com/learn/latex/Learn_LaTeX_in_30_minutes
\end{document}
